\section{Fundamental Theorem for Injective PEPS}\label{sec:peps_ft}

The edge construction is transported from single edges to arbitrary regions.
Reading a region $R$ and its complement $V\setminus R$ as two-block tensors,
when the blocked tensors of $B$ on $R$ and $V\setminus R$ are injective and the
bonds are positive, the region-insertion identities determine the boundary
conjugation uniquely: two invertible matrices $Z_1,Z_2$ realizing the same
insertion identity satisfy, for every boundary matrix $M$,
\[
    Z_1\,M\,Z_1^{-1}=Z_2\,M\,Z_2^{-1}.
\]
Comparing one vertex $v$ against its complement
then upgrades this to $A_v=\lambda_v\widetilde B_v$ at each vertex, once the edge
gauges have been absorbed into $\widetilde B$, and the per-edge gauges combine
into one gauge relating the two tensors. This proves the Fundamental Theorem for
injective PEPS, Theorem~\ref{thm:fundamentalTheorem_PEPS}, the formalization of
\cite[Theorem~2]{Molnar2018NormalPEPS}.

\begin{definition}[Region two-block tensor]%
    \label{def:peps_regionTwoBlock}
    \lean{TNLean.PEPS.regionTwoBlock}
    \leanok
    For a finite region $R\subseteq V$, the blocked tensor of $R$ may be read
    as a two-block tensor with a one-point external boundary and shared bonds
    the edges crossing the boundary of $R$:
    \[
        B_{A,R}(*,\eta,\sigma)=T_{A,R}^{\eta}(\sigma).
    \]
    This is the region block used in the one-block-against-complement step of
    \cite[Section~3]{Molnar2018NormalPEPS}.
\end{definition}

\begin{definition}[Complement-boundary identification]%
    \label{def:peps_regionComplementBoundaryConfig}
    \lean{TNLean.PEPS.regionBoundaryEdgeToCompl,
        TNLean.PEPS.regionBoundaryEdgeComplEquiv,
        TNLean.PEPS.regionComplementBoundaryConfig}
    \leanok
    The edges crossing the boundary of $R$ are the same edges crossing the
    boundary of $V\setminus R$. Thus a boundary configuration $\eta$ for $R$
    determines the boundary configuration $\eta^c$ for the complement by reading
    the same edge labels through this identification.
\end{definition}

\begin{definition}[Complement-region two-block tensor]%
    \label{def:peps_regionComplementTwoBlock}
    \lean{TNLean.PEPS.regionComplementTwoBlock}
    \leanok
    \uses{def:peps_regionComplementBoundaryConfig}
    For a finite region $R\subseteq V$, the complement block $V\setminus R$ is
    a two-block tensor over the same crossing bonds:
    \[
        B_{A,R^c}(*,\eta,\tau)
        =
        T_{A,V\setminus R}^{\eta^c}(\tau).
    \]
    This is the complementary block in the region/complement comparison of
    \cite[Section~3]{Molnar2018NormalPEPS}.
\end{definition}

\begin{theorem}[Injectivity of the complement-region two-block tensor]%
    \label{thm:peps_isTwoBlockInjective_regionComplementTwoBlock}
    \lean{TNLean.PEPS.isTwoBlockInjective_regionComplementTwoBlock}
    \leanok
    \uses{def:peps_regionComplementTwoBlock, def:peps_regionInjectivityData}
    If the blocked tensor family of $V\setminus R$ is injective, then the
    complement-region two-block tensor is injective:
    \[
        \operatorname{inj}(T_{A,V\setminus R})
        \Longrightarrow
        \operatorname{inj}(B_{A,R^c}).
    \]
\end{theorem}
\begin{proof}\leanok
    The one-point external boundary does not change the coefficient family, and
    the complement-boundary identification only reindexes the crossing bonds.
    More explicitly, if $\eta_1,\eta_2$ are boundary configurations on
    $\partial R$, then
    \[
        \eta_1^c=\eta_2^c\quad\Longrightarrow\quad \eta_1=\eta_2 .
    \]
    Thus the coefficients of $T_{A,V\setminus R}$ and $B_{A,R^c}$ differ only
    by an injective relabelling of boundary configurations.
\end{proof}

\begin{definition}[Region insertion coefficient]%
    \label{def:peps_regionInsertedCoeff}
    \lean{TNLean.PEPS.regionInsertedCoeff}
    \leanok
    \uses{def:peps_regionTwoBlock, def:peps_regionComplementTwoBlock}
    Let $f$ be an edge crossing the boundary of $R$, and let
    $M\in\MN{D_f}$.  The region insertion coefficient is
    \[
        C_A(R,f,M;\sigma,\tau)
        =
        \sum_{\substack{\mu,\nu\\
        \mu|_{\partial R\setminus\{f\}}
        =\nu|_{\partial R\setminus\{f\}}}}
        M_{\mu_f,\nu_f}\,
        T_{A,R}^{\mu}(\sigma)\,
        T_{A,V\setminus R}^{\nu^c}(\tau).
    \]
    It is the coefficient obtained by inserting $M$ on the boundary bond $f$
    between the region $R$ and its complement.
\end{definition}

\begin{theorem}[Region insertion coefficients determine the inserted matrix]%
    \label{thm:peps_regionInsertedCoeff_injective}
    \lean{TNLean.PEPS.regionInsertedCoeff_injective}
    \leanok
    \uses{def:peps_regionInsertedCoeff}
    Suppose that the blocked tensor maps of $A$ on $R$ and on
    $V\setminus R$ are injective, and that every bond dimension of $A$ is
    positive.  If $f\in\partial R$ and
    $M,M'\in\MN{D_A(f)}$ satisfy
    \[
        C_A(R,f,M;\sigma,\tau)
        =
        C_A(R,f,M';\sigma,\tau)
    \]
    for every physical configuration $\sigma$ on $R$ and $\tau$ on
    $V\setminus R$, then $M=M'$.
\end{theorem}
\begin{proof}
    \leanok
    Subtracting the two coefficients gives
    \[
        C_A(R,f,M-M';\sigma,\tau)=0
    \]
    for every $\sigma,\tau$.  Injectivity of the region and complement
    blocked tensor maps reads off every entry of $M-M'$ from configurations
    agreeing away from $f$, so $M-M'=0$ and hence $M=M'$.
\end{proof}

\begin{definition}[Bond-inserted region insert]%
    \label{def:peps_bond_inserted_region_insert}
    \lean{TNLean.PEPS.bondInsertedRegionInsert}
    \leanok
    Let $f$ be an edge crossing the boundary of $R$, and let
    $M\in\MN{D_f}$.  The bond-inserted region insert
    $I_{A,R,f,M}$ is the insert on $R$ defined by
    \[
        (I_{A,R,f,M})_{\nu,\sigma}
        =
        \sum_{\substack{\mu\in\partial R\\
        \mu|_{\partial R\setminus\{f\}}
        =\nu|_{\partial R\setminus\{f\}}}}
        M_{\mu_f,\nu_f}\,T_{A,R}^{\mu}(\sigma).
    \]
    Thus the complement boundary value $\nu$ is fixed away from $f$, while
    the matrix $M$ couples the two $f$-labels.  This is the one-bond insertion
    in the region/complement coefficient of
    \cite[Section~3]{Molnar2018NormalPEPS}.
\end{definition}

\begin{theorem}[Bond insertion as a deformed region state]
    \label{thm:peps_deformedRegionState_bondInsertedRegionInsert}
    \lean{TNLean.PEPS.deformedRegionState_bondInsertedRegionInsert}
    \leanok
    \uses{def:peps_bond_inserted_region_insert,
        def:peps_regionInsertedCoeff}
    For every physical configuration $\sigma$ on $R$ and $\tau$ on
    $V\setminus R$, the deformed state of the bond-inserted region insert is
    the region insertion coefficient:
    \[
        \Psi_{A,R,I_{A,R,f,M}}(\sigma,\tau)
        =
        C_A(R,f,M;\sigma,\tau).
    \]
\end{theorem}
\begin{proof}
    \leanok
    Expanding the deformed state gives
    \[
        \Psi_{A,R,I_{A,R,f,M}}(\sigma,\tau)
        =
        \sum_{\nu\in\partial R}
        I_{A,R,f,M}(\nu,\sigma)\,
        T_{A,V\setminus R}^{\nu^c}(\tau).
    \]
    Substituting the definition of $I_{A,R,f,M}$ and exchanging the two finite
    sums gives
    \[
    \begin{aligned}
        \Psi_{A,R,I_{A,R,f,M}}(\sigma,\tau)
        &=
        \sum_{\nu\in\partial R}
        \left(
        \sum_{\substack{\mu\in\partial R\\
        \mu|_{\partial R\setminus\{f\}}
        =\nu|_{\partial R\setminus\{f\}}}}
        M_{\mu_f,\nu_f}\,T_{A,R}^{\mu}(\sigma)
        \right)
        T_{A,V\setminus R}^{\nu^c}(\tau)\\
        &=
        C_A(R,f,M;\sigma,\tau).
    \end{aligned}
    \]
\end{proof}

\begin{theorem}[Region insertion as an edge insertion]
    \label{thm:peps_regionInsertedCoeff_eq_smul_edgeInsertedCoeff}
    \lean{TNLean.PEPS.regionInsertedCoeff_eq_smul_edgeInsertedCoeff}
    \leanok
    \uses{def:peps_regionInsertedCoeff, def:peps_edgeInsertedCoeff}
    Let $f$ be a boundary edge of $R$, let $N\in\MN{D_B(f)}$, let $\sigma$
    be a physical configuration on $R$, and let $\tau$ be a physical
    configuration on $V\setminus R$.  Put
    \[
        p_B(R)=\prod_{g\notin\partial R}D_B(g).
    \]
    Then the region insertion coefficient factors through the edge-inserted
    coefficient:
    \[
        C_B(R,f,N;\sigma,\tau)
        =
        p_B(R)\,
        E_B\bigl(f,\operatorname{Assemble}_R(\sigma,\tau),
            \mathcal O_{R,f}(N)\bigr).
    \]
    Here $\mathcal O_{R,f}$ is the boundary-edge orientation, equal to the
    identity if the left endpoint of $f$ lies in $R$ and to transpose
    otherwise.
\end{theorem}
\begin{proof}
    \leanok
    Let $\mathcal C_{R,f}(\sigma,\tau)$ denote the open-edge configurations
    compatible with the two physical configurations and with the two boundary
    readings agreeing away from $f$.  For
    $\xi\in\mathcal C_{R,f}(\sigma,\tau)$, let
    $\operatorname{Fib}_{R,f}(\xi)$ be the corresponding fibre of agreeing
    boundary readings.  Grouping the expansion of $C_B(R,f,N;\sigma,\tau)$ by
    such an open-edge configuration $\xi$ gives
    \[
    \begin{aligned}
        C_B(R,f,N;\sigma,\tau)
        &=
        \sum_{\xi\in\mathcal C_{R,f}(\sigma,\tau)}
        \#\operatorname{Fib}_{R,f}(\xi)\,
        N_{\xi_{R,f}^{(1)},\xi_{R,f}^{(2)}}
        \prod_{v\in V}
        B_v\bigl(\xi_v,\operatorname{Assemble}_R(\sigma,\tau)_v\bigr),\\
        \#\operatorname{Fib}_{R,f}(\xi)
        &=
        \prod_{g\notin\partial R}D_B(g)
        =
        p_B(R).
    \end{aligned}
    \]
    Substituting this fibre cardinality gives
    \[
    \begin{aligned}
        C_B(R,f,N;\sigma,\tau)
        &=
        p_B(R)
        \sum_{\xi\in\mathcal C_{R,f}(\sigma,\tau)}
        N_{\xi_{R,f}^{(1)},\xi_{R,f}^{(2)}}
        \prod_{v\in V}
        B_v\bigl(\xi_v,\operatorname{Assemble}_R(\sigma,\tau)_v\bigr)\\
        &=
        p_B(R)\,
        E_B\bigl(f,\operatorname{Assemble}_R(\sigma,\tau),
            \mathcal O_{R,f}(N)\bigr).
    \end{aligned}
    \]
    The second equality is the definition of the edge-inserted coefficient,
    with the inserted matrix read in the boundary-edge orientation.
\end{proof}

\begin{theorem}[Bond insertion through a corner extension]
    \label{thm:peps_regionInsertedCoeff_eq_extendInsert_bondInserted}
    \lean{TNLean.PEPS.regionInsertedCoeff_eq_extendInsert_bondInserted}
    \leanok
    \uses{thm:peps_deformedRegionState_bondInsertedRegionInsert,
        thm:peps_corner_extension_state_preserves}
    Let $R\subseteq S$, and suppose that all bond dimensions are positive.
    For every global physical configuration $\sigma$,
    \[
        C_A\bigl(R,f,M;\sigma|_R,\sigma|_{V\setminus R}\bigr)
        =
        \Psi_{A,S,
        \operatorname{Ext}_{R\subseteq S}(I_{A,R,f,M})}(\sigma).
    \]
    In particular, the one-bond insertion on $R$ may be read after extending
    the corresponding insert to any larger region $S$.
\end{theorem}
\begin{proof}
    \leanok
    The preceding theorem and state preservation for corner extension give
    \[
    \begin{aligned}
        C_A\bigl(R,f,M;\sigma|_R,\sigma|_{V\setminus R}\bigr)
        &=
        \Psi_{A,R,I_{A,R,f,M}}
            (\sigma|_R,\sigma|_{V\setminus R})\\
        &=
        \Psi_{A,S,
            \operatorname{Ext}_{R\subseteq S}(I_{A,R,f,M})}(\sigma).
    \end{aligned}
    \]
\end{proof}

\begin{theorem}[Region insertion as a two-block insertion]%
    \label{thm:peps_twoBlockInsertedCoeff_eq_regionInsertedCoeff}
    \lean{TNLean.PEPS.twoBlockInsertedCoeff_eq_regionInsertedCoeff}
    \leanok
    \uses{def:peps_regionInsertedCoeff, def:peps_regionTwoBlock,
        def:peps_regionComplementTwoBlock}
    The abstract two-block insertion coefficient of the pair
    $(B_{A,R},B_{A,R^c})$ is the region insertion coefficient:
    \[
        C_{B_{A,R},B_{A,R^c}}(f,M;*,*,\sigma,\tau)
        =
        C_A(R,f,M;\sigma,\tau).
    \]
\end{theorem}
\begin{proof}\leanok
    Expanding the abstract two-block coefficient gives exactly the two boundary
    sums in the displayed formula for $C_A(R,f,M;\sigma,\tau)$.
\end{proof}

\begin{theorem}[Region insertion under bond-dimension transport]%
    \label{thm:peps_regionInsertedCoeff_reindexTensor}
    \lean{TNLean.PEPS.regionInsertedCoeff_reindexTensor}
    \leanok
    \uses{def:peps_regionInsertedCoeff}
    If a tensor $B$ is transported along an equality of bond dimensions, then
    the region insertion coefficient is transported by the induced
    matrix-algebra identification on the inserted bond:
    \[
        C_{B^h}(R,f,N;\sigma,\tau)
        =
        C_B(R,f,\phi_f(N);\sigma,\tau).
    \]
\end{theorem}
\begin{proof}\leanok
    Reindex both boundary-configuration sums by the bond-dimension
    equivalences. The constraint away from $f$ and the two blocked-region
    weights are preserved, while the inserted matrix entry becomes the
    transported entry $\phi_f(N)$.
\end{proof}

\begin{theorem}[Same two-block insertions from region insertions]%
    \label{thm:peps_sameTwoBlockInsertions_of_regionInsertedCoeff_eq}
    \lean{TNLean.PEPS.sameTwoBlockInsertions_of_regionInsertedCoeff_eq}
    \leanok
    \uses{thm:peps_twoBlockInsertedCoeff_eq_regionInsertedCoeff,
        thm:peps_regionInsertedCoeff_reindexTensor}
    Suppose $A$ and $B$ have identified bond dimensions. Write $B^h$ for the
    transported tensor, and write $\phi_f$ for the induced transport on the
    boundary bond $f$. If
    \[
        C_A(R,f,N;\sigma,\tau)
        =
        C_B(R,f,\phi_f(N);\sigma,\tau)
    \]
    for every boundary edge $f$, matrix $N$, and physical configurations
    $\sigma,\tau$, then the two pairs
    \[
        (B_{A,R},B_{A,R^c}),\qquad
        (B_{B^h,R},B_{B^h,R^c})
    \]
    have the same one-bond insertion coefficients after the same
    bond-dimension transport.
\end{theorem}
\begin{proof}\leanok
    Rewrite the two abstract insertion coefficients as region insertion
    coefficients and use the bond-dimension transport formula:
    \[
        C_{B_{A,R},B_{A,R^c}}(f,N;*,*,\sigma,\tau)
        = C_A(R,f,N;\sigma,\tau)
        = C_B(R,f,\phi_f(N);\sigma,\tau)
        = C_{B_{B^h,R},B_{B^h,R^c}}(f,N;*,*,\sigma,\tau).
    \]
\end{proof}

\begin{theorem}[Swapping the two blocks in a one-bond insertion]
    \label{thm:peps_twoBlockInsertedCoeff_swap}
    \lean{TNLean.PEPS.twoBlockInsertedCoeff_swap}
    \leanok
    Let \(B_1,B_2\) be two blocks joined along a shared bond \(f\).  Then inserting
    a matrix \(X\) while contracting \(B_1\) against \(B_2\) is the same as
    contracting \(B_2\) against \(B_1\) with the transposed insertion:
    \[
        C_{B_1,B_2}(f,X;\eta_1,\eta_2,\sigma_1,\sigma_2)
        =
        C_{B_2,B_1}(f,X^T;\eta_2,\eta_1,\sigma_2,\sigma_1).
    \]
\end{theorem}
\begin{proof}\leanok
    Exchange the two boundary-configuration sums.  The condition that the two
    configurations agree away from \(f\) is symmetric, and
    \(X^T_{\nu_f,\mu_f}=X_{\mu_f,\nu_f}\).
\end{proof}

\begin{theorem}[Complement-side reading of a region insertion]
    \label{thm:peps_regionInsertedCoeff_complement_reading}
    \lean{TNLean.PEPS.regionInsertedCoeff_eq_complementTwoBlock}
    \leanok
    \uses{thm:peps_twoBlockInsertedCoeff_swap,
        thm:peps_twoBlockInsertedCoeff_eq_regionInsertedCoeff}
    For a boundary edge \(f\in\partial R\), the region-inserted coefficient can be
    read by contracting the complement block first:
    \[
        C_A(R,f,X;\sigma,\tau)
        =
        C_{A,R^c;R}(f,X^T;\tau,\sigma).
    \]
    Here \(C_{A,R^c;R}\) denotes the same one-bond insertion coefficient with the
    two blocks \(R^c\) and \(R\) interchanged.
\end{theorem}
\begin{proof}\leanok
    Identify \(C_A(R,f,X;\sigma,\tau)\) with the two-block insertion of
    \((B_{A,R},B_{A,R^c})\), then apply
    Theorem~\ref{thm:peps_twoBlockInsertedCoeff_swap}.
\end{proof}

\begin{theorem}[Incident form gives the transferred endpoint form]
    \label{thm:peps_hform_of_region_incident_form}
    \lean{TNLean.PEPS.hform_of_isIncidentMatrixForm}
    \leanok
    \uses{thm:peps_physical_to_virtual_insertion}
    Assume that \(A\)'s and \(B\)'s components at the in-region endpoint of \(f\)
    are linearly independent, and that \(D_B(e)>0\) for every edge \(e\).
    Let \(W\) be the virtual pullback, at the in-region endpoint of \(f\), of the
    physical operator obtained by inserting \(X^T\) on \(f\).  If
    \[
        W=\mathcal I_f(P)
    \]
    is in incident-matrix form on the boundary leg \(f\), then the transfer matrix
    read off from \(W\) reinserts to \(W\).
\end{theorem}
\begin{proof}\leanok
    The transfer matrix is defined by reading the incident matrix out of \(W\).
    Reading \(\mathcal I_f(P)\) gives \(P\), and reinserting \(P\) returns \(W\).
\end{proof}

\begin{theorem}[Virtual pullback realizes the endpoint operator]
    \label{thm:peps_region_insertion_virtual_pullback_realizes}
    \lean{TNLean.PEPS.localTensorMap_localVirtualOpOfPhysicalOpAt_regionInsertionOp}
    \leanok
    \uses{def:peps_same_state, def:IsVertexInjective, def:peps_localTensorMap,
        thm:peps_physical_to_virtual_insertion}
    Assume that \(A\) and \(B\) have the same PEPS state, are vertex-injective,
    and have positive bond dimensions.  Let \(v\) be the in-region endpoint of
    \(f\), and let \(W\) be the virtual pullback, for \(B\) at \(v\), of the
    physical operator obtained by inserting \(X^T\) on \(f\).  Then for every
    local virtual coefficient \(c\),
    \[
        \mathcal T_{B,v}(Wc)
        =
        O_{A,R,f}(X^T)\,\mathcal T_{B,v}(c).
    \]
\end{theorem}
\begin{proof}\leanok
    The physical endpoint operator preserves the image of
    \(\mathcal T_{B,v}\). The chosen left inverse defining the virtual pullback
    therefore realizes the same action after applying \(\mathcal T_{B,v}\).
\end{proof}

\begin{definition}[Out-of-region endpoint and blocked-region inverse]
    \label{def:peps_region_out_endpoint_block_inverse}
    \lean{TNLean.PEPS.regionBoundaryEdgeOutVertex,
        TNLean.PEPS.regionBoundaryEdgeOutVertex_not_mem,
        TNLean.PEPS.regionBoundaryEdgeOutVertex_mem_compl,
        TNLean.PEPS.regionBoundaryEdgeInVertex_compl_eq_outVertex,
        TNLean.PEPS.regionBlockedTensorMap,
        TNLean.PEPS.regionBlockedTensorMap_injective_of_injective,
        TNLean.PEPS.regionBlockedTensorMap_ker_eq_bot_of_injective,
        TNLean.PEPS.regionBlockedLeftInverse,
        TNLean.PEPS.regionBlockedLeftInverse_comp_regionBlockedTensorMap,
        TNLean.PEPS.regionBlockedLeftInverse_apply_regionBlockedTensorMap,
        TNLean.PEPS.regionBlockedTensorMap_apply,
        TNLean.PEPS.regionBlockedTensorMap_single,
        TNLean.PEPS.regionBlockedLeftInverse_regionBlockedWeight}
    \leanok
    \uses{def:peps_regionInjectivityData}
    The out-of-region endpoint of \(f\in\partial R\) is the endpoint outside
    \(R\).  It is the in-region endpoint of the same edge viewed as an edge of
    \(\partial(V\setminus R)\).  The blocked-region map is
    \[
        c\longmapsto \sum_{\mu\in\partial R} c_\mu\,T_{A,R}^{\mu},
    \]
    and injectivity of the blocked tensor gives a left inverse.  In particular,
    \[
        L_{A,R}\!\left(T_{A,R}^{\mu}\right)=e_\mu .
    \]
\end{definition}

\begin{theorem}[Complement-side cast and endpoint realization]
    \label{thm:peps_regionInsertedCoeff_complement_cast}
    \lean{TNLean.PEPS.mem_compl_compl_iff,
        TNLean.PEPS.regionDoubleComplVertexEquiv,
        TNLean.PEPS.regionDoubleComplPhysicalConfig,
        TNLean.PEPS.isRegionBoundaryEdge_compl_compl_iff,
        TNLean.PEPS.regionBoundaryEdgeDoubleComplEquiv,
        TNLean.PEPS.regionDoubleComplBoundaryConfig,
        TNLean.PEPS.regionBlockedWeight_doubleCompl,
        TNLean.PEPS.regionComplementBoundaryConfigEquiv,
        TNLean.PEPS.regionComplementBoundaryConfig_apply_toCompl,
        TNLean.PEPS.regionComplementBoundaryConfig_compl_compl,
        TNLean.PEPS.sameAwayFromBond_complementBoundaryConfig,
        TNLean.PEPS.regionInsertedCoeff_eq_compl,
        TNLean.PEPS.regionInsertedCoeff_eq_regionRealizationSum_compl,
        TNLean.PEPS.regionInsertedCoeff_eq_smul_op_regionStateVec_compl}
    \leanok
    \uses{def:peps_region_out_endpoint_block_inverse,
        thm:peps_regionInsertedCoeff_complement_reading}
    With \(f'\) the boundary edge \(f\) read on \(V\setminus R\), one has
    \[
        C_A(R,f,X;\sigma,\tau)
        =
        C_A(V\setminus R,f',X^T;\tau,\sigma^{cc}),
    \]
    where \(\sigma^{cc}\) is \(\sigma\) transported across
    \(V\setminus(V\setminus R)=R\).  Consequently the same coefficient is realized
    through the endpoint of \(f\) lying outside \(R\).
\end{theorem}
\begin{proof}\leanok
    Reindex the two boundary-configuration sums by the boundary-edge equivalence
    between \(R\) and its complement.  The double complement carries the blocked
    weight of \(V\setminus(V\setminus R)\) back to the blocked weight of \(R\).
    Applying the usual endpoint realization to \(V\setminus R\) gives the displayed
    out-of-region endpoint realization.
\end{proof}

\begin{theorem}[Region resonate identity and blocked read-offs]
    \label{thm:peps_region_resonate_identity_readoffs}
    \lean{TNLean.PEPS.regionInsertionOp_regionStateVec_pin_compl,
        TNLean.PEPS.region_resonate_identity,
        TNLean.PEPS.regionComplementRow,
        TNLean.PEPS.regionInsertedCoeff_eq_complement_blockedMap,
        TNLean.PEPS.regionRegionRow,
        TNLean.PEPS.regionInsertedCoeff_eq_region_blockedMap,
        TNLean.PEPS.regionBlockedLeftInverse_complement_regionInsertedCoeff,
        TNLean.PEPS.regionBlockedLeftInverse_region_regionInsertedCoeff}
    \leanok
    \uses{def:peps_region_out_endpoint_block_inverse,
        thm:peps_regionInsertedCoeff_complement_cast}
    Assume that \(A\)'s components at the in-region endpoint of \(f\) and at the
    in-complement-region endpoint of the complementary boundary edge are linearly
    independent, and that \(A\) and \(B\) define the same PEPS state.  Let
    \(P_R\) and \(P_{R^c}\) be the interior bond products of \(R\) and \(R^c\).
    Let \(O^{\mathrm{in}}_f(X)\) and \(O^{\mathrm{out}}_f(X)\) be the physical
    endpoint operators obtained from the two virtual insertions, and let
    \(\Psi_B(R;\sigma,\tau)\) and \(\Psi_B(R^c;\tau,\sigma^{cc})\) be the
    corresponding closed \(B\)-state vectors.  Then the two endpoint readings of
    the same inserted coefficient give
    \[
        P_R\,O^{\mathrm{in}}_f(X)\,\Psi_B(R;\sigma,\tau)
        =
        P_{R^c}\,O^{\mathrm{out}}_f(X)\,\Psi_B(R^c;\tau,\sigma^{cc}).
    \]
    With \(\rho^{R^c}_{X,\sigma}\) and \(\rho^{R}_{X,\tau}\) denoting the two
    row functions, one has
    \[
        C_A(R,f,X;\sigma,-)
        =
        T_{A,R^c}\!\left(\rho^{R^c}_{X,\sigma}\right),
        \qquad
        C_A(R,f,X;-,\tau)
        =
        T_{A,R}\!\left(\rho^{R}_{X,\tau}\right).
    \]
    Hence the blocked left inverses read off the rows:
    \[
        L_{A,R^c}\!\left(C_A(R,f,X;\sigma,-)\right)
        =
        \rho^{R^c}_{X,\sigma},
        \qquad
        L_{A,R}\!\left(C_A(R,f,X;-,\tau)\right)
        =
        \rho^{R}_{X,\tau}.
    \]
\end{theorem}
\begin{proof}\leanok
    The in-region and out-of-region endpoint pins both evaluate to
    \(C_A(R,f,X;\sigma,\tau)\), hence they agree.  Expanding
    \(C_A(R,f,X;\sigma,\tau)\) as a sum over boundary configurations isolates,
    respectively, the \(R^c\)-blocked and \(R\)-blocked tensor maps; the left
    inverses then read off the row functions.
\end{proof}

\begin{definition}[Incident form of the region resonance]
    \label{def:peps_region_resonate_reconcile}
    \lean{TNLean.PEPS.RegionResonateReconcile}
    \leanok
    \uses{thm:peps_region_resonate_identity_readoffs}
    The incident-form condition associated with the region resonance asserts that,
    for every matrix \(X\) inserted on \(f\), the virtual pullback \(W_X\) of the
    in-region endpoint operator is in incident-matrix form:
    \[
        \exists Y,\qquad W_X=\mathcal I_f(Y).
    \]
\end{definition}

\begin{definition}[Region insertion transfer]
    \label{def:peps_region_insertion_transfer}
    \lean{TNLean.PEPS.RegionInsertionTransfer}
    \leanok
    \uses{def:peps_regionInsertedCoeff}
    A region insertion transfer along a boundary edge $f\in\partial R$ consists
    of maps
    \[
        T_f:\MN{D_A(f)}\to\MN{D_B(f)},
        \qquad
        S_f:\MN{D_B(f)}\to\MN{D_A(f)}
    \]
    such that, for every inserted matrix and every pair of physical
    configurations,
    \[
        C_A(R,f,X;\sigma,\tau)
        =
        C_B(R,f,T_f(X);\sigma,\tau),
        \qquad
        C_B(R,f,Y;\sigma,\tau)
        =
        C_A(R,f,S_f(Y);\sigma,\tau).
    \]
    The forward map is required to satisfy
    \[
        T_f(XY)=T_f(X)T_f(Y),\qquad T_f(I)=I.
    \]
\end{definition}

\begin{definition}[Region physical-to-virtual realization]
    \label{def:peps_region_transfer_realizes}
    \lean{TNLean.PEPS.RegionTransferRealizes}
    \leanok
    \uses{thm:peps_region_insertion_virtual_pullback_realizes}
    Let $v$ be the endpoint of $f$ lying in $R$.  A matrix transfer $T_f$ is
    realized on the region if, for every matrix $X$ inserted on $A$'s boundary
    bond and every local virtual coefficient $c$ for $B$ at $v$,
    \[
        O_{A,R,f}(X^T)\,\mathcal T_{B,v}(c)
        =
        \mathcal T_{B,v}\!\left(\mathcal I_f(T_f(X)^T)c\right).
    \]
    This is the region version of the physical-to-virtual insertion relation in
    \cite[Section~3, lines~254--582]{Molnar2018NormalPEPS}.
\end{definition}

\begin{theorem}[Realized region transfers are multiplicative and unital]
    \label{thm:peps_region_transfer_realized_maps}
    \lean{TNLean.PEPS.regionTransferMatrix_mul_of_realizes,
        TNLean.PEPS.regionTransferMatrix_one_of_realizes}
    \leanok
    \uses{def:peps_region_transfer_realizes, def:peps_region_insertion_transfer}
    If a region matrix transfer is realized in the sense of
    Definition~\ref{def:peps_region_transfer_realizes}, then its forward map is
    multiplicative:
    \[
        T_f(XY)=T_f(X)T_f(Y).
    \]
    If in addition the two PEPS tensors have the same state and the same bond
    dimensions, all bond dimensions of $B$ are positive, and the blocked
    tensors of $B$ on $R$ and its complement are injective, then the identity
    insertion gives
    \[
        T_f(I)=I.
    \]
\end{theorem}
\begin{proof}\leanok
    For multiplicativity, the physical insertion for $(XY)^T$ is the composite
    of the physical insertions for $Y^T$ and $X^T$.  Realizing these insertions
    on the $B$-side gives two incident-matrix actions on the same local tensor
    image:
    \[
        \mathcal I_f(T_f(XY)^T)
        =
        \mathcal I_f((T_f(X)T_f(Y))^T).
    \]
    Reading the incident matrix on the boundary leg gives
    $T_f(XY)=T_f(X)T_f(Y)$.  For the identity, the coefficient
    $C_A(R,f,I;\sigma,\tau)$ is the interior bond product times the state
    coefficient of $A$, hence equals $C_B(R,f,I;\sigma,\tau)$ by equality of
    states and bond dimensions.  The two blocked-injectivity hypotheses and
    positivity give injectivity of the $B$-side region-inserted coefficient,
    whence $T_f(I)=I$.
\end{proof}

\begin{definition}[Region insertion transfer from realized maps]
    \label{def:peps_region_insertion_transfer_of_realizes}
    \lean{TNLean.PEPS.regionInsertionTransfer_of_realizes}
    \leanok
    \uses{def:peps_region_transfer_realizes,
        thm:peps_region_transfer_realized_maps}
    Suppose the forward and reverse boundary matrix transfers are realized on
    the two tensors, the two tensors have the same state and the same bond
    dimensions, and the region and complement-region blocked maps of $B$ are
    injective.
    Then the maps $T_f$ and $S_f$ form a region insertion transfer:
    \[
        C_A(R,f,X;\sigma,\tau)
        =
        C_B(R,f,T_f(X);\sigma,\tau),
        \qquad
        C_B(R,f,Y;\sigma,\tau)
        =
        C_A(R,f,S_f(Y);\sigma,\tau),
    \]
    together with $T_f(XY)=T_f(X)T_f(Y)$ and $T_f(I)=I$.
\end{definition}

\begin{theorem}[Incident form gives the region insertion transfer and gauge]
    \label{thm:peps_region_reconcile_to_transfer_gauge}
    \lean{TNLean.PEPS.regionTransferRealizes_of_isIncidentMatrixForm,
        TNLean.PEPS.regionTransferRealizes_of_reconcile,
        TNLean.PEPS.regionInsertionTransfer_of_reconcile,
        TNLean.PEPS.exists_regionEdgeGauge_of_reconcile}
    \leanok
    \uses{def:peps_region_resonate_reconcile,
        thm:peps_hform_of_region_incident_form, def:peps_same_state,
        def:IsVertexInjective, def:peps_regionInjectivityData}
    Assume that \(A\) and \(B\) have the same PEPS state, are vertex-injective,
    have positive bond dimensions, and have injective blocked maps for \(R\) and
    \(R^c\). If the incident-form condition holds in both directions, then the
    matrix reading gives maps
    \[
        T_f:\MN{D_A(f)}\to\MN{D_B(f)},
        \qquad
        S_f:\MN{D_B(f)}\to\MN{D_A(f)}
    \]
    realizing all region insertion coefficients. If, in addition,
    \(D_A(e)=D_B(e)\) for every edge \(e\), then there is an invertible gauge
    \(Z_f\in\GL(D_B(f))\) such that
    \[
        T_f(X)=Z_f\,\iota_f(X)\,Z_f^{-1},
    \]
    where \(\iota_f:\MN{D_A(f)}\simeq\MN{D_B(f)}\) is the identification induced
    by the equality of bond dimensions.
\end{theorem}
\begin{proof}\leanok
    The incident-form condition supplies incident-matrix form for every inserted
    matrix. Theorem~\ref{thm:peps_hform_of_region_incident_form} gives the
    realization identities. The forward and reverse realization identities
    assemble the transfer datum, and the one-bond algebra theorem reads it as
    conjugation by an invertible matrix.
\end{proof}

\begin{theorem}[Coefficient transfer gives incident form]
    \label{thm:peps_region_coeff_transfer_to_reconcile}
    \lean{TNLean.PEPS.regionInsertionOp_regionStateVec_pin_leg,
        TNLean.PEPS.regionInsertionOp_realizes_of_realizes_on_stateVec,
        TNLean.PEPS.regionInteriorBondProd_pos,
        TNLean.PEPS.regionInsertionOp_regionStateVec_eq_of_coeff_eq,
        TNLean.PEPS.isIncidentMatrixForm_of_coeff_eq,
        TNLean.PEPS.regionResonateReconcile_of_coeff_transfer}
    \leanok
    \uses{def:peps_region_resonate_reconcile}
    Assume that \(A\) and \(B\) define the same PEPS state, that \(B\) is
    vertex-injective, that \(D_B(e)>0\) for every edge \(e\), and that
    \(D_A(e)=D_B(e)\) for every edge \(e\).  Assume also that \(A\)'s and \(B\)'s
    components at the in-region endpoint of \(f\) are linearly independent.
    Suppose that for
    every matrix \(X\) on \(A\)'s boundary bond there is a matrix \(Y\) on \(B\)'s
    boundary bond such that
    \[
        C_A(R,f,X;\sigma,\tau)=C_B(R,f,Y;\sigma,\tau)
    \]
    for all \(\sigma,\tau\).  Then the virtual pullback of \(X\) is in
    incident-matrix form on \(f\).
\end{theorem}
\begin{proof}\leanok
    Equality of the coefficients gives equality on the closed state vectors, one
    physical leg at a time.  These state-vector coefficients span the local
    virtual coefficient space, so the equality extends to all coefficients.  The
    virtual pullback is therefore the incident insertion determined by \(Y\).
\end{proof}

\begin{theorem}[First coefficient through the second complement block]
    \label{thm:peps_region_first_coeff_second_complement_block}
    \lean{TNLean.PEPS.regionInteriorBondProd_smul_regionStateVec,
        TNLean.PEPS.regionBlockedWeight_update_eq_regionWeightVec,
        TNLean.PEPS.regionInteriorBondProd_smul_regionStateVec_eq_sum,
        TNLean.PEPS.regionInsertedCoeff_eq_complement_blockedMap_B,
        TNLean.PEPS.regionBlockedLeftInverse_complement_regionInsertedCoeff_B}
    \leanok
    \uses{def:peps_region_out_endpoint_block_inverse}
    Assume that \(B\)'s complement blocked tensor map is injective, that \(A\)
    and \(B\) define the same PEPS state, that \(D_A(e)=D_B(e)\) for every edge
    \(e\), and that \(A\)'s component at the in-region endpoint of \(f\) is
    linearly independent.
    The coefficient \(C_A(R,f,X;\sigma,\tau)\), as a function of the complement
    physical configuration \(\tau\), lies in the image of \(B\)'s complement
    blocked map.  Applying the complement blocked left inverse of \(B\) recovers
    the row obtained by applying \(A\)'s endpoint operator to \(B\)'s region weight
    vectors.
\end{theorem}
\begin{proof}\leanok
    The identity insertion expresses the closed state vector as a contraction of
    the region and complement blocked weights.  Substituting the endpoint operator
    into this expression gives the complement-blocked factorization, and the left
    inverse reads off the corresponding row.
\end{proof}

\begin{theorem}[Transfer coefficient and double factorization]
    \label{thm:peps_region_transfer_coeff_double_factorization}
    \lean{TNLean.PEPS.vSideRow,
        TNLean.PEPS.regionInsertedCoeff_eq_complement_blockedMap_vSideRow,
        TNLean.PEPS.regionInsertedCoeff_eq_of_complementRow_eq,
        TNLean.PEPS.regionDoubleComplBoundaryConfigEquiv,
        TNLean.PEPS.regionBlockedTensorMap_doubleCompl,
        TNLean.PEPS.complSideRow,
        TNLean.PEPS.regionInsertedCoeff_eq_region_blockedMap_B,
        TNLean.PEPS.regionRowB,
        TNLean.PEPS.regionRowB_eq_complSideRow,
        TNLean.PEPS.regionRowB_mem_range,
        TNLean.PEPS.transferCoeff,
        TNLean.PEPS.regionRowB_eq_complement_blockedMap,
        TNLean.PEPS.regionInsertedCoeff_eq_doubleSum_transferCoeff,
        TNLean.PEPS.regionInsertedCoeff_eq_of_transferCoeff_form,
        TNLean.PEPS.vSideRow_eq_region_blockedMap_transferCoeff}
    \leanok
    \uses{thm:peps_region_first_coeff_second_complement_block}
    Assume that \(B\)'s region and complement blocked tensor maps are injective,
    that \(A\)'s components at the in-region endpoint of \(f\) and at the
    in-complement-region endpoint of the complementary boundary edge are linearly
    independent, that \(A\) and \(B\) define the same PEPS state, and that
    \(D_A(e)=D_B(e)\) for every edge \(e\).
    For fixed \(X\) and \(\sigma\), define the row
    \[
        v_X^\sigma(\nu)
        =
        \left[
            O_{A,R,f}(X^T)\,
            T_{B,R}^{\tilde\nu}(\sigma)
            \right]_{\sigma(v_f)},
    \]
    where \(\tilde\nu\) is the \(R\)-boundary configuration corresponding to the
    \(R^c\)-boundary configuration \(\nu\), and \(v_f\) is the endpoint of \(f\)
    in \(R\).
    The first tensor's inserted coefficient can be written as a double contraction
    of the second tensor's region and complement blocked weights:
    \[
        C_A(R,f,X;\sigma,\tau)
        =
        \sum_{\mu,\nu}
        K_X(\mu,\nu)\,T_{B,R}^{\mu}(\sigma)\,
        T_{B,R^c}^{\nu}(\tau).
    \]
    If \(K_X\) has the incident form
    \[
        K_X(\mu,\nu)=\mathbf 1_{\mu|_{\partial R\setminus\{f\}}
        =\nu|_{\partial R\setminus\{f\}}}\,Y_{\mu_f,\nu_f},
    \]
    then \(C_A(R,f,X;\sigma,\tau)=C_B(R,f,Y;\sigma,\tau)\).
\end{theorem}
\begin{proof}\leanok
    First factor \(C_A\) through \(B\)'s complement block, then through \(B\)'s
    region block.  The two left inverses define the coefficient \(K_X\).  Substituting
    the displayed incident form of \(K_X\) turns the double contraction into the
    defining double sum for \(C_B(R,f,Y;\sigma,\tau)\).
\end{proof}

\begin{theorem}[Incident form of the transfer coefficient]
    \label{thm:peps_region_transfer_coeff_incident_form}
    \lean{TNLean.PEPS.transferCoeff_column_eq_regionBlockedLeftInverse_vSideRow,
        TNLean.PEPS.vSideRow_eq_localTensorMap_virtualPullback,
        TNLean.PEPS.vSideRow_eq_incidentSum_of_virtualPullback_incidentForm,
        TNLean.PEPS.transferCoeff_eq_incidentForm_of_virtualPullback_incidentForm}
    \leanok
    \uses{thm:peps_region_transfer_coeff_double_factorization,
        thm:peps_hform_of_region_incident_form,
        thm:peps_region_insertion_virtual_pullback_realizes}
    Under the hypotheses of
    Theorem~\ref{thm:peps_region_transfer_coeff_double_factorization}, assume in
    addition that \(A\) and \(B\) are vertex-injective, that \(D_A(e)>0\) and
    \(D_B(e)>0\) for every edge \(e\), and that \(B\)'s component at the
    in-region endpoint of \(f\) is linearly independent.
    If the virtual pullback \(W_X\) is in incident-matrix form
    \(W_X=\mathcal I_f(Y^T)\), then the transfer coefficient \(K_X\) has the
    incident form determined by \(Y\):
    \[
        K_X(\mu,\nu)=\mathbf 1_{\mu|_{\partial R\setminus\{f\}}
        =\nu|_{\partial R\setminus\{f\}}}\,Y_{\mu_f,\nu_f}.
    \]
\end{theorem}
\begin{proof}\leanok
    The column of \(K_X\) at fixed \(\nu\) is the region blocked left inverse of
    the row \(v_X^\sigma\).  If \(W_X=\mathcal I_f(Y^T)\), this row is precisely the
    incident-matrix sum of \(Y\) against the region blocked weights.  The left
    inverse sends each blocked weight to the corresponding standard boundary
    configuration, giving the displayed formula.
\end{proof}

\begin{definition}[Coefficient transfer map]
    \label{def:peps_coeff_transfer_map}
    \lean{TNLean.PEPS.coeffTransferMap}
    \leanok
    \uses{def:peps_regionInsertedCoeff}
    Suppose that every matrix $X\in\MN{D_A(f)}$ has a matrix
    $Y\in\MN{D_B(f)}$ with matching region-inserted coefficients:
    \[
        C_A(R,f,X;\sigma,\tau)
        =
        C_B(R,f,Y;\sigma,\tau).
    \]
    Choosing one such $Y$ for each $X$ defines a forward coefficient transfer
    $T_f:\MN{D_A(f)}\to\MN{D_B(f)}$.
\end{definition}

\begin{theorem}[Coefficient transfer matches coefficients and preserves the identity]
    \label{thm:peps_coeff_transfer_map_properties}
    \lean{TNLean.PEPS.coeffTransferMap_coeff,
        TNLean.PEPS.coeffTransferMap_one}
    \leanok
    \uses{def:peps_coeff_transfer_map}
    The chosen coefficient transfer satisfies, for every inserted matrix $X$ and
    every pair of physical configurations,
    \[
        C_A(R,f,X;\sigma,\tau)
        =
        C_B(R,f,T_f(X);\sigma,\tau).
    \]
    If the tensors have the same state and the same bond dimensions, all
    $B$-bond dimensions are positive, and the region and complement-region
    blocked maps of $B$ are injective, then
    \[
        T_f(I)=I.
    \]
\end{theorem}
\begin{proof}\leanok
    The coefficient equality is the defining property of the chosen transfer.
    For $X=I$, both region-inserted coefficients are the corresponding interior
    bond products times the common PEPS state coefficient:
    \[
        C_A(R,f,I;\sigma,\tau)
        =
        C_B(R,f,I;\sigma,\tau).
    \]
    The two blocked-map injectivity hypotheses and positivity of the $B$-bond
    dimensions make the $B$-side region-inserted coefficient injective.  Hence
    the chosen matrix $T_f(I)$ is equal to $I$.
\end{proof}

\begin{definition}[Region insertion transfer from coefficient transfers]
    \label{def:peps_region_insertion_transfer_of_coeff_transfer}
    \lean{TNLean.PEPS.regionInsertionTransfer_of_coeffTransfer}
    \leanok
    \uses{def:peps_region_insertion_transfer,
        def:peps_coeff_transfer_map,
        thm:peps_coeff_transfer_map_properties}
    Suppose coefficient transfers exist in both directions:
    \[
        C_A(R,f,X;\sigma,\tau)
        =
        C_B(R,f,T_f(X);\sigma,\tau),
        \qquad
        C_B(R,f,Y;\sigma,\tau)
        =
        C_A(R,f,S_f(Y);\sigma,\tau).
    \]
    Assume also that the two tensors have the same state and the same bond
    dimensions, all $B$-bond dimensions are positive, and the region and
    complement-region blocked maps of $B$ are injective.  If the forward choice
    is multiplicative,
    \[
        T_f(XY)=T_f(X)T_f(Y),
    \]
    then $T_f$ and $S_f$ form a region insertion transfer.
\end{definition}

\begin{theorem}[Per-edge gauge from coefficient transfers]
    \label{thm:peps_region_edge_gauge_of_coeff_transfer}
    \lean{TNLean.PEPS.exists_regionEdgeGauge_of_coeffTransfer}
    \leanok
    \uses{def:peps_region_insertion_transfer_of_coeff_transfer}
    Assume coefficient transfers exist in both directions, the forward transfer
    is multiplicative, the two tensors have the same state and the same bond
    dimensions, all bond dimensions are positive, and the region and
    complement-region blocked maps for both tensors are injective.  Then there
    is an invertible boundary gauge $Z_f$ such that, for every
    $X\in\MN{D_A(f)}$,
    \[
        T_f(X)=Z_f\,\iota_f(X)\,Z_f^{-1},
    \]
    where $\iota_f:\MN{D_A(f)}\simeq\MN{D_B(f)}$ is the identification induced by
    the equality of bond dimensions.
\end{theorem}
\begin{proof}\leanok
    The coefficient transfers form a region insertion transfer by
    Definition~\ref{def:peps_region_insertion_transfer_of_coeff_transfer}.  The
    insertion-transfer algebra isomorphism gives an algebra equivalence between
    the two boundary matrix algebras, and the Skolem--Noether theorem for matrix
    algebras writes this equivalence as conjugation by an invertible matrix.
\end{proof}

\begin{theorem}[Determinacy of a region insertion transfer]
    \label{thm:peps_regionInsertedCoeff_transferMap_unique}
    \lean{TNLean.PEPS.regionInsertedCoeff_transferMap_unique}
    \leanok
    \uses{def:peps_regionInsertedCoeff, def:peps_region_out_endpoint_block_inverse}
    Assume that the blocked tensors of \(B\) on \(R\) and \(R^c\) are injective
    and that all virtual bonds of \(B\) have positive dimension.  Let
    \(T_1,T_2:\MN{D_A(f)}\to\MN{D_B(f)}\) be two maps.  For \(i=1,2\), suppose
    \(T_i\) satisfies, for every inserted matrix and every pair of physical
    configurations,
    \[
        C_A(R,f,M;\sigma,\tau)
        =
        C_B(R,f,T_i(M);\sigma,\tau).
    \]
    Then \(T_1(M)=T_2(M)\) for every \(M\).
\end{theorem}
\begin{proof}\leanok
    For every pair of physical configurations, the two displayed identities give
    \[
        C_B(R,f,T_1(M);\sigma,\tau)
        =
        C_A(R,f,M;\sigma,\tau)
        =
        C_B(R,f,T_2(M);\sigma,\tau).
    \]
    Injectivity of the region and complement blocked tensors, together with
    positivity of all virtual bonds of \(B\), gives
    \[
        \left(
        \forall\sigma,\tau,\,
        C_B(R,f,T_1(M);\sigma,\tau)
        =
        C_B(R,f,T_2(M);\sigma,\tau)
        \right)
        \Longrightarrow
        T_1(M)=T_2(M).
    \]
\end{proof}

\begin{theorem}[Reindexing commutes with conjugation]
    \label{thm:peps_reindexAlgEquiv_gaugeConj}
    \lean{TNLean.PEPS.reindexAlgEquiv_gaugeConj}
    \leanok
    Assume \(D=E\), let \(\iota:\MN{D}\simeq\MN{E}\) be the canonical
    identification of matrix algebras, and let \(Z\in\GL(D)\).  Then for every
    \(N\in\MN{D}\),
    \[
        \iota(ZNZ^{-1})
        =
        \iota(Z)\,\iota(N)\,\iota(Z)^{-1}.
    \]
\end{theorem}
\begin{proof}\leanok
    The map \(\iota\) is an algebra equivalence, so it preserves products and
    inverses.  Applying this to the three factors \(Z\), \(N\), and \(Z^{-1}\)
    gives the displayed identity.
\end{proof}

\begin{theorem}[Coefficient identities determine the boundary conjugation]
    \label{thm:peps_gaugeConj_eq_of_coeffIdentities}
    \lean{TNLean.PEPS.gaugeConj_eq_of_coeffIdentities}
    \leanok
    \uses{thm:peps_regionInsertedCoeff_transferMap_unique}
    Assume that the blocked tensors of \(B\) on \(R\) and \(R^c\) are injective
    and that all virtual bonds of \(B\) have positive dimension.  Suppose that
    the equality \(D_A(f)=D_B(f)\) fixes the canonical identification
    \[
        \iota:\MN{D_A(f)}\simeq\MN{D_B(f)}.
    \]
    For \(i=1,2\), suppose that the invertible matrix \(Z_i\in\GL(D_B(f))\)
    realizes the region-insertion identity
    \[
        C_A(R,f,M;\sigma,\tau)
        =
        C_B(R,f,Z_i\,\iota(M)\,Z_i^{-1};\sigma,\tau).
    \]
    Then the two conjugations agree on every boundary matrix:
    \[
        Z_1\,\iota(M)\,Z_1^{-1}
        =
        Z_2\,\iota(M)\,Z_2^{-1}.
    \]
\end{theorem}
\begin{proof}\leanok
    For \(i=1,2\), define the transfer map
    \[
        T_i(M):=Z_i\,\iota(M)\,Z_i^{-1}.
    \]
    Its coefficient identity is
    \[
        C_A(R,f,M;\sigma,\tau)
        =
        C_B(R,f,T_i(M);\sigma,\tau).
    \]
    The preceding determinacy theorem gives
    \[
        T_1(M)=T_2(M).
    \]
    Substituting the definitions of \(T_1\) and \(T_2\) yields
    \[
        Z_1\,\iota(M)\,Z_1^{-1}
        =
        Z_2\,\iota(M)\,Z_2^{-1}.
    \]
\end{proof}

%======================================================================
% InsertSplit: the one-site quotient of the blocked-region weight
%======================================================================

\begin{theorem}[Region vertex product across an inserted site]%
    \label{thm:peps_prod_region_insert_split}
    \lean{TNLean.PEPS.prod_region_insert_split}
    \leanok
    Let $v\notin R$.  At a fixed global virtual configuration $\zeta$ and a
    physical configuration $\sigma$ on $R\cup\{v\}$, the product of the local
    tensors over the vertices of $R\cup\{v\}$ factors as the local tensor at the
    inserted site $v$ times the product over the vertices of $R$:
    \[
        \prod_{w\in R\cup\{v\}}A_w\bigl(\zeta|_{\operatorname{inc}(w)},\sigma_w\bigr)
        =
        A_v\bigl(\zeta|_{\operatorname{inc}(v)},\sigma_v\bigr)
        \prod_{w\in R}A_w\bigl(\zeta|_{\operatorname{inc}(w)},\sigma|_R{}_w\bigr).
    \]
\end{theorem}
\begin{proof}\leanok
    The inserted site $v$ is isolated from the product over $R\cup\{v\}$ by
    splitting off a single factor, and the residual product over the remaining
    vertices is reindexed through the bijection between the vertices of $R$ and
    the non-$v$ vertices of $R\cup\{v\}$.
\end{proof}

\begin{theorem}[One-site split of the blocked-region weight]%
    \label{thm:peps_regionBlockedWeight_insert_eq_sum_split}
    \lean{TNLean.PEPS.regionBlockedWeight_insert_eq_sum_split}
    \leanok
    \uses{thm:peps_prod_region_insert_split}
    Let $v\notin R$.  The blocked-region weight of $R\cup\{v\}$ is the constrained
    sum, over global virtual configurations $\zeta$ whose crossing-edge label on
    $R\cup\{v\}$ is $\mu$, of the inserted-site tensor times the vertex product
    over $R$:
    \[
        T_{A,R\cup\{v\}}^{\mu}(\sigma)
        =
        \sum_{\zeta:\,\zeta|_{\partial(R\cup\{v\})}=\mu}
        A_v\bigl(\zeta|_{\operatorname{inc}(v)},\sigma_v\bigr)
        \prod_{w\in R}A_w\bigl(\zeta|_{\operatorname{inc}(w)},\sigma|_R{}_w\bigr).
    \]
\end{theorem}
\begin{proof}\leanok
    Each summand of the blocked-region weight is the vertex product over
    $R\cup\{v\}$, which factors by
    Theorem~\ref{thm:peps_prod_region_insert_split}.
\end{proof}

\begin{theorem}[A non-incident crossing edge survives the insertion]%
    \label{thm:peps_isRegionBoundaryEdge_insert_of_not_incident}
    \lean{TNLean.PEPS.isRegionBoundaryEdge_insert_of_not_incident}
    \leanok
    Let $v\notin R$.  A crossing edge of $R$ whose two endpoints both differ from
    $v$ is also a crossing edge of $R\cup\{v\}$.
\end{theorem}
\begin{proof}\leanok
    A crossing edge of $R$ has exactly one endpoint inside $R$.  Adjoining $v$ to
    $R$ leaves both endpoints on the same side as before, since neither endpoint
    is $v$, so the edge still crosses the boundary of $R\cup\{v\}$.
\end{proof}

\begin{definition}[Boundary label of the smaller region across an insertion]%
    \label{def:peps_boundaryLabelOfInsert}
    \lean{TNLean.PEPS.boundaryLabelOfInsert}
    \leanok
    \uses{thm:peps_isRegionBoundaryEdge_insert_of_not_incident}
    Let $v\notin R$.  Given a crossing-edge label $\mu$ for $R\cup\{v\}$ and a
    local virtual configuration $\eta$ on the edges incident to $v$, the
    crossing-edge label of $R$ is read off these two data: a crossing edge of $R$
    incident to $v$ has become internal to $R\cup\{v\}$, so its label is taken
    from $\eta$; a crossing edge of $R$ not incident to $v$ is also a crossing
    edge of $R\cup\{v\}$, so its label is taken from $\mu$.
\end{definition}

\begin{theorem}[Boundary-label fiber identity across an inserted site]%
    \label{thm:peps_regionBoundaryLabel_eq_boundaryLabelOfInsert}
    \lean{TNLean.PEPS.regionBoundaryLabel_eq_boundaryLabelOfInsert}
    \leanok
    \uses{def:peps_boundaryLabelOfInsert}
    Let $v\notin R$.  If a global virtual configuration $\zeta$ restricts to $\mu$
    on the crossing edges of $R\cup\{v\}$ and to $\eta$ on the edges incident to
    $v$, then the crossing-edge label of $R$ read from $\zeta$ is exactly the
    bridge label of $\mu$ and $\eta$.
\end{theorem}
\begin{proof}\leanok
    The two labels are compared edge by edge.  On a crossing edge of $R$
    incident to $v$ both sides read $\zeta$ on an edge incident to $v$, which is
    $\eta$.  On a crossing edge of $R$ not incident to $v$ both sides read
    $\zeta$ on a crossing edge of $R\cup\{v\}$, which is $\mu$.
\end{proof}

%======================================================================
% Gauge-absorption bridge at region granularity
%======================================================================

\begin{theorem}[Region gauge absorption]%
    \label{thm:peps_regionInsertedCoeff_applyGauge}
    \lean{TNLean.PEPS.regionInsertedCoeff_applyGauge}
    \leanok
    \uses{thm:peps_regionInsertedCoeff_eq_smul_edgeInsertedCoeff,
        thm:peps_edge_gauge_absorption, def:applyGauge}
    Let $X$ be an edge-gauge family, write $B^X$ for the gauged tensor, and let
    $f$ be a crossing edge of $R$.  Inserting $M$ on $f$ between the region and
    its complement in $B^X$ equals inserting, in $B$, the matrix obtained by
    orienting $M$ to the boundary-edge convention, conjugating by the transpose
    of the endpoint gauge $X_f$, and orienting back:
    \[
        C_{B^X}(R,f,M;\sigma,\tau)
        =
        C_B\bigl(R,f,
            \mathcal O_{R,f}\bigl(X_f^{\mathsf T}\,
                \mathcal O_{R,f}(M)\,(X_f^{-1})^{\mathsf T}\bigr);
            \sigma,\tau\bigr),
    \]
    where $\mathcal O_{R,f}$ is the boundary-edge orientation.  This is the
    region-granularity analogue of
    Theorem~\ref{thm:peps_edge_gauge_absorption}.
\end{theorem}
\begin{proof}\leanok
    By Theorem~\ref{thm:peps_regionInsertedCoeff_eq_smul_edgeInsertedCoeff} the
    region/complement contraction is the single-bond cut at $f$ up to the
    interior-bond multiplicity, which is unchanged by a gauge.  At the edge
    granularity every interior edge and every crossing edge other than $f$
    cancels its gauge pairwise, while the two endpoint gauges on $f$ conjugate
    the inserted matrix.  Transporting the cancelled identity back through the
    region-to-edge identity, and using that the boundary-edge orientation is an
    involution, gives the displayed equality.
\end{proof}

%======================================================================
% Region-transport covariance and torus translation covariance
%======================================================================

\begin{theorem}[Transport of the complement-side blocked weight]%
    \label{thm:peps_regionComplementWeight_transport}
    \lean{TNLean.PEPS.regionComplementWeight_transport}
    \leanok
    \uses{def:peps_tensor_transport, def:peps_regionComplementTwoBlock}
    Let $\phi:G\simeq G'$ be a graph isomorphism.  The blocked weight of the
    transported tensor $A^\phi$ over the complement of the image region, read at
    the image complement boundary configuration and the image complement physical
    configuration, equals the blocked weight of $A$ over $V\setminus R$:
    \[
        T_{A^\phi,V'\setminus\phi(R)}^{(\phi\cdot\eta)^c}(\phi\cdot\tau)
        =
        T_{A,V\setminus R}^{\eta^c}(\tau).
    \]
\end{theorem}
\begin{proof}\leanok
    The image of the complement of $R$ is the complement of the image of $R$, and
    after this identification the blocked weight transports along $\phi$ by the
    change of summation variables on virtual configurations.
\end{proof}

\begin{theorem}[Transport preserves agreement away from a bond]%
    \label{thm:peps_sameAwayFromBond_regionBoundaryConfigMap}
    \lean{TNLean.PEPS.sameAwayFromBond_regionBoundaryConfigMap}
    \leanok
    \uses{thm:peps_regionComplementWeight_transport}
    Let $\phi:G\simeq G'$ be a graph isomorphism and let $f$ be a crossing edge
    of $R$.  Two crossing-edge labels of $R$ agree away from $f$ if and only if
    their images under the boundary reindex induced by $\phi$ agree away from the
    image of $f$.
\end{theorem}
\begin{proof}\leanok
    The boundary reindex is a bijection on crossing edges that carries $f$ to its
    image, so it carries the crossing edges other than $f$ bijectively to the
    crossing edges other than the image of $f$.  Agreement on the former family
    is therefore equivalent to agreement on the latter.
\end{proof}

\begin{theorem}[Covariance of the region insertion coefficient]%
    \label{thm:peps_regionInsertedCoeff_transport}
    \lean{TNLean.PEPS.regionInsertedCoeff_transport}
    \leanok
    \uses{def:peps_regionInsertedCoeff, def:peps_tensor_transport,
        thm:peps_regionComplementWeight_transport,
        thm:peps_sameAwayFromBond_regionBoundaryConfigMap}
    Let $\phi:G\simeq G'$ be a graph isomorphism and let $f$ be a crossing edge
    of $R$.  The region insertion coefficient of the transported tensor $A^\phi$
    over the image region, with $M$ inserted on the image of $f$ and the physical
    configurations relabelled by $\phi$, equals the region insertion coefficient
    of $A$ over $R$, with $M$ carried back to the bond of $A$ at $f$:
    \[
        C_{A^\phi}\bigl(\phi(R),\phi(f),M;\phi\cdot\sigma,\phi\cdot\tau\bigr)
        =
        C_A\bigl(R,f,\phi^{*}M;\sigma,\tau\bigr).
    \]
    No endpoint orientation enters: the matrix is inserted between the two single
    boundary values on the edge.
\end{theorem}
\begin{proof}\leanok
    The two crossing-label sums reindex by the boundary action of $\phi$, the
    agreement-away-from-$f$ constraint is preserved by
    Theorem~\ref{thm:peps_sameAwayFromBond_regionBoundaryConfigMap}, the region
    weight transports along $\phi$, and the complement weight transports by
    Theorem~\ref{thm:peps_regionComplementWeight_transport}.  The inserted matrix
    entry is read at the two boundary values, which the bond-dimension
    identification carries to the corresponding entry of $\phi^{*}M$.
\end{proof}

\begin{theorem}[Region insertion coefficient under an equality of tensors]%
    \label{thm:peps_regionInsertedCoeff_congrTensor}
    \lean{TNLean.PEPS.regionInsertedCoeff_congrTensor}
    \leanok
    \uses{def:peps_regionInsertedCoeff}
    If two tensors are equal, then their region insertion coefficients are equal,
    with the inserted matrix carried across the induced equality of bond
    dimensions:
    \[
        A_1=A_2
        \quad\Longrightarrow\quad
        C_{A_1}(R,f,M;\sigma,\tau)
        =
        C_{A_2}\bigl(R,f,\psi M;\sigma,\tau\bigr).
    \]
    This bridges the literal tensor of a translation-invariant PEPS with its
    transported copy.
\end{theorem}
\begin{proof}\leanok
    Substituting the equality of tensors identifies the two coefficients, and the
    induced equality of bond dimensions is the identity, so the carried matrix is
    the original one.
\end{proof}

\begin{theorem}[Bond dimension at a translated crossing edge]%
    \label{thm:peps_bondDim_boundaryEdgeMap_translate}
    \lean{TNLean.PEPS.bondDim_boundaryEdgeMap_translate}
    \leanok
    \uses{def:peps_torus_translation_invariant,
        def:peps_torus_geometry_translation}
    Let $A$ be translation invariant on the discrete torus, let $f$ be a crossing
    edge of $R$, and let $\tau_{a,b}$ be a translation.  The bond dimension of $A$
    at the translated crossing edge equals its bond dimension at the reference
    crossing edge:
    \[
        D_A\bigl(\tau_{a,b}(f)\bigr)=D_A(f).
    \]
\end{theorem}
\begin{proof}\leanok
    Translation carries the reference edge to its image, and translation
    invariance fixes the tensor under transport, so the bond dimension is
    unchanged.
\end{proof}

\begin{theorem}[Translation covariance of the coefficient identity]%
    \label{thm:peps_regionInsertedCoeff_translate_coeffIdentity}
    \lean{TNLean.PEPS.regionInsertedCoeff_translate_coeffIdentity}
    \leanok
    \uses{thm:peps_regionInsertedCoeff_transport,
        thm:peps_regionInsertedCoeff_congrTensor,
        thm:peps_bondDim_boundaryEdgeMap_translate,
        def:peps_torus_translation_invariant}
    Let $A$ and $B$ be translation invariant on the discrete torus, and let $f$ be
    a crossing edge of $R$.  Suppose a map $g$ realizes the coefficient identity
    at the reference edge:
    \[
        C_A(R,f,M;\sigma,\tau)
        =
        C_B\bigl(R,f,g(M);\sigma,\tau\bigr).
    \]
    Then, for every translation $\tau_{a,b}$, the same identity holds at the
    translated crossing edge of the translated region, with the inserted matrix
    carried between the translated and reference bonds:
    \[
        C_A\bigl(\tau_{a,b}(R),\tau_{a,b}(f),M';
            \tau_{a,b}\cdot\sigma,\tau_{a,b}\cdot\tau\bigr)
        =
        C_B\bigl(\tau_{a,b}(R),\tau_{a,b}(f),
            \widehat{g}(M');
            \tau_{a,b}\cdot\sigma,\tau_{a,b}\cdot\tau\bigr),
    \]
    where $\widehat{g}$ applies $g$ to the reference-bond reindex of $M'$ and
    carries the result back to the translated bond.
\end{theorem}
\begin{proof}\leanok
    Translation invariance identifies the transported tensor with the original,
    so Theorem~\ref{thm:peps_regionInsertedCoeff_transport} carries the reference
    coefficient of $A$ to the translated edge with no change of tensor;
    Theorem~\ref{thm:peps_regionInsertedCoeff_congrTensor} mediates the literal
    and transported tensors.  Applying the reference identity for $g$ and
    transporting back for $B$ in the same way gives the displayed equality, the
    inserted matrix being carried across the equal bond dimensions of
    Theorem~\ref{thm:peps_bondDim_boundaryEdgeMap_translate}.
\end{proof}

\begin{definition}[Single-vertex two-block tensor]%
    \label{def:peps_vertexTwoBlock}
    \lean{TNLean.PEPS.vertexTwoBlock}
    \leanok
    \uses{def:peps_tensor}
    Fix a vertex $v$.  The single-vertex tensor $A_v$ is regarded as a
    two-block tensor whose shared bonds are the edges incident to $v$, whose
    external boundary is a one-point space, and whose physical index is the
    physical index at $v$:
    \[
        T_v(*,\eta,\sigma)=A_v(\eta,\sigma).
    \]
\end{definition}

\begin{theorem}[Injectivity of the single-vertex two-block tensor]%
    \label{thm:peps_isTwoBlockInjective_vertexTwoBlock}
    \lean{TNLean.PEPS.isTwoBlockInjective_vertexTwoBlock}
    \leanok
    \uses{def:IsVertexInjective, def:peps_vertexTwoBlock}
    If $A$ is vertex-injective, then the single-vertex two-block tensor at every
    vertex is injective.
\end{theorem}
\begin{proof}\leanok
    Vertex injectivity is linear independence of the coefficient family
    $\eta\mapsto A_v(\eta,-)$.  The auxiliary one-point boundary only reindexes
    this same family, so the two-block injectivity condition is the same linear
    independence statement.
\end{proof}

\begin{definition}[Complement two-block tensor]%
    \label{def:peps_complementTwoBlock}
    \lean{TNLean.PEPS.complementTwoBlock}
    \leanok
    \uses{def:peps_tensor}
    For a vertex $v$, the complementary block $V\setminus\{v\}$ defines a
    two-block tensor on the same incident shared bonds. Its coefficient is the
    contraction of all tensors away from $v$, with the incident bond indices
    left open at the complement endpoints:
    \[
        T_{v^c}(*,\eta,\tau)=A_{v^c}(\eta,\tau).
    \]
\end{definition}

\begin{theorem}[Injectivity of the vertex-complement contraction]%
    \label{thm:peps_vertexComplementTensorInjective}
    \lean{TNLean.PEPS.vertexComplementTensorInjective_of_isVertexInjective}
    \leanok
    \uses{def:IsVertexInjective, def:peps_tensor}
    If $A$ is vertex-injective and every virtual bond space has positive
    dimension, then the coefficient family
    \[
        \eta\longmapsto \bigl(\tau\longmapsto A_{v^c}(\eta,\tau)\bigr)
    \]
    of the vertex-complement contraction is linearly independent.
\end{theorem}
\begin{proof}\leanok
    Let $c=(c_\eta)_\eta$ be a finite coefficient family on the open star of
    $v$. For a set $S\subseteq V\setminus\{v\}$, let $K_S(c)$ be the assertion
    that for every virtual configuration $\zeta_0$ and physical configuration
    $\tau$,
    \[
        \sum_{\zeta}
        \mathbf 1[
                \zeta_f=\zeta_{0,f}\text{ for every edge }f
                \text{ whose endpoints are outside }S]\,
        c_{\zeta|_{\operatorname{Star}(v)}}
        \prod_{w\in S} A_w(\zeta|_{\operatorname{Star}(w)},\tau_w)=0 .
    \]
    The initial vanishing of the complement contraction is exactly
    $K_{V\setminus\{v\}}(c)$. If $j\in S$ and $j\neq v$, the one-sided inverse at
    the vertex $j$ separates the local tensor $A_j$ and gives
    $K_S(c)\Rightarrow K_{S\setminus\{j\}}(c)$. Iterating this implication over
    the complement vertices gives $K_\varnothing(c)$. At $S=\varnothing$ the
    exposed virtual indices are all fixed by $\zeta_0$; positive bond dimensions
    allow $\zeta_0$ to realize any prescribed star configuration at $v$. Hence
    every coefficient $c_\eta$ is zero.
\end{proof}

\begin{theorem}[Injectivity of the complement two-block tensor]%
    \label{thm:peps_isTwoBlockInjective_complementTwoBlock}
    \lean{TNLean.PEPS.isTwoBlockInjective_complementTwoBlock}
    \leanok
    \uses{def:IsVertexInjective, def:peps_complementTwoBlock,
        thm:peps_vertexComplementTensorInjective}
    If $A$ is vertex-injective and every virtual bond space has positive
    dimension, then the complement two-block tensor at every vertex is
    injective.
\end{theorem}
\begin{proof}\leanok
    \uses{thm:peps_vertexComplementTensorInjective}
    The previous theorem gives linear independence of the coefficient family
    $\eta\mapsto A_{v^c}(\eta,-)$. The two-block tensor only adds a one-point
    external boundary, so this is the same injectivity condition after the
    harmless reindexing $(*,\eta)\leftrightarrow\eta$.
\end{proof}

\begin{theorem}[One-vertex versus complement comparison]%
    \label{thm:peps_oneVertexComplementComparison}
    \lean{TNLean.PEPS.one_vertex_complement_comparison}
    \leanok
    \uses{thm:peps_postAbsorptionEdgeInsertionEquality,
        thm:peps_twoInjectiveTensorInsertionComparison}
    Let the shared bond type and the two external boundary spaces be nonempty.
    After the edge gauges have been absorbed into the second tensor family
    $\widetilde B$, block one vertex $v$ against its complement. Assume that
    $A_v,A_{v^c},\widetilde B_v,\widetilde B_{v^c}$ are injective as two-block
    tensors. The post-absorption insertion equality gives, for every shared
    bond $b$, every matrix $X\in\MN{D_b}$, and all external and
    physical indices,
    \[
        C_{A_v,A_{v^c}}(b,X;\eta_v,\eta_{v^c},\sigma_v,\sigma_{v^c})
        =
        C_{\widetilde B_v,\widetilde B_{v^c}}
        (b,X;\eta_v,\eta_{v^c},\sigma_v,\sigma_{v^c}).
    \]
    Then
    \[
        \exists\lambda_v\in\C^\times,\qquad
        A_v=\lambda_v\widetilde B_v.
    \]
    \begin{center}
        \TNPEPSOneVertexComplementComparison
    \end{center}
    This is the final local comparison following the post-absorption insertion
    equality in \cite[Section~3]{Molnar2018NormalPEPS}.
\end{theorem}
\begin{proof}\leanok
    With $A_1=A_v$, $A_2=A_{v^c}$, $B_1=\widetilde B_v$, and
    $B_2=\widetilde B_{v^c}$, the post-absorption equality is precisely
    \[
        \forall b,X,\eta_v,\eta_{v^c},\sigma_v,\sigma_{v^c},\quad
        C_{A_1,A_2}(b,X;\eta_v,\eta_{v^c},\sigma_v,\sigma_{v^c})
        =
        C_{B_1,B_2}(b,X;\eta_v,\eta_{v^c},\sigma_v,\sigma_{v^c}).
    \]
    Theorem~\ref{thm:peps_twoInjectiveTensorInsertionComparison} gives
    \[
        \exists\lambda_v\in\mathbb C^\times,\qquad
        A_v=\lambda_v\widetilde B_v,\qquad
        A_{v^c}=\lambda_v^{-1}\widetilde B_{v^c}.
    \]
    The desired local proportionality is the first equality.
\end{proof}

\begin{remark}[Normal-region form]
    The normal-region application in \cite[Section~4]{Molnar2018NormalPEPS}
    uses the same scalar comparison after comparing two injective regions
    $R\subset S=R\cup\{v\}$. If the region comparisons give
    $A_R=\lambda_R\widetilde B_R$ and
    $A_S=\lambda_S\widetilde B_S$, then the source-paper diagram reads
    \[
        A_R A_v
        =
        A_S
        =
        \lambda_S\widetilde B_S
        =
        \lambda_S\widetilde B_R\widetilde B_v.
    \]
    Applying the left inverse supplied by injectivity of the block
    $A_R=\lambda_R\widetilde B_R$ gives
    $A_v=(\lambda_S/\lambda_R)\widetilde B_v$.
\end{remark}

\begin{theorem}[Fundamental Theorem for injective PEPS, conditional on
        bond-dimension equality]%
    \label{thm:fundamentalTheorem_PEPS_of_bondDim}
    \lean{TNLean.PEPS.fundamentalTheorem_PEPS_of_bondDim}
    \leanok
    \uses{thm:gaugeConsistency, thm:peps_oneVertexComplementComparison,
        def:GaugeEquivPEPS, def:IsVertexInjective}
    Suppose that the corresponding virtual edge spaces of $A$ and $B$ are
    already identified, every virtual bond of $A$ has positive dimension, and
    the graph is connected.
    If $A$ and $B$ are vertex-injective and give the same PEPS state, then
    there are invertible matrices $X_e$, one for each edge, such that at every
    vertex the tensor $B_v$ is obtained from $A_v$ by the oriented endpoint
    action of the incident matrices.
\end{theorem}
\begin{proof}\leanok
    Gauge consistency supplies the edge gauges directly, and the gauge
    equivalence packages them with the identified bond dimensions.
\end{proof}

\begin{theorem}[Fundamental Theorem for injective PEPS]%
    \label{thm:fundamentalTheorem_PEPS}
    \lean{TNLean.PEPS.fundamentalTheorem_PEPS}
    \leanok
    % Scope restriction: this is the positive-bond variant of Theorem 2 in
    % \cite{Molnar2018NormalPEPS}. The elimination plan for the positivity
    % hypotheses is recorded in
    % \path{docs/paper-gaps/peps_injective_ft_section3_route.tex}.
    \uses{thm:fundamentalTheorem_PEPS_of_bondDim, def:GaugeEquivPEPS,
        def:IsVertexInjective}
    Let $A$ and $B$ be vertex-injective PEPS tensors on a finite simple
    graph, and suppose every virtual bond of $A$ and $B$ has positive
    dimension and that the graph is connected. If they give the same PEPS
    state, then their corresponding
    virtual edge spaces have the same dimensions and, after identifying these
    spaces, there are invertible matrices $X_e$, one for each edge, such that
    $B_v$ is obtained from $A_v$ by the oriented endpoint action of the matrices
    incident to $v$.
\end{theorem}
\begin{proof}\leanok
    The bond dimensions agree edgewise, because blocking around an edge gives
    two injective three-site chains whose matched matrix insertions on that bond
    form an algebra equivalence between the two bond matrix algebras, forcing
    equal sizes. With the bond dimensions identified, the conditional form of
    the theorem applies.
\end{proof}
